\section{Fixture Pattern}

The intended fixture family follows the four-role W-series pattern. These are
defined roles for later implementation, not a claim that fixture files or tests
already exist.

\begin{center}
\small
\begin{tabularx}{\textwidth}{lX}
\toprule
Fixture & Role \\
\midrule
No-anomaly & A synthetic count table drawn from a plausible election-count
distribution with no planted digit manipulation. A cautious method should return
low priority or explain ordinary sampling variation. \\
Injected tampered digit & A table with a planted digit-pattern distortion, such
as excess repeated final digits or altered leading digits in a subset. The
method should surface the affected unit while still labeling the result as
investigative. \\
Innocent digit deviation & A bounded precinct-size table that deviates from a
Benford-style reference for structural reasons. This exercises the false-positive
warning and should not be described as suspicious by itself. \\
Boundary failure & A table that formally ``fails'' a digit test but is fully
explained by precinct-size design, turnout ceilings, or pooling choices. The
report must carry the explanation in \texttt{investigation\_note}. \\
\bottomrule
\end{tabularx}
\end{center}

The last two fixtures are as important as the planted one. A digit method that
only finds the injected distortion is incomplete; it must also demonstrate that
ordinary administrative structure can produce high statistics and that the
report language keeps those statistics in their proper, weak-signal role.
